\lecture{13}
\title{Neural Networks: Depth and Optimization}

\begin{document}

\subsection{Neural Network Depth}

In this section,
we provide an example of a \emph{deep} ReLU-activated network
that cannot be accurately represented
by a \emph{shallow} ReLU-activated network
(without using an excessive number of nodes).
Formally, we show:

\begin{theorem}{Depth Necessity}
    For every $L \in \mathbb{N}$,
    we can construct a function $f \colon [0, 1] \to [0, 1]$
    such that:
    \begin{itemize}
        \item It is exactly representable
        by a ReLU network of depth $O(L^2)$
        with $O(L^2)$ nodes.
        \item For any ReLU network $h \colon [0, 1] \to \mathbb{R}$
        of depth $L$ with at most $2^L$ nodes,
        we have $\int_{0}^1 |f(x) - h(x)| \, \mathrm{d}x \geq \frac{1}{8}$.
    \end{itemize}
\end{theorem}

\begin{remark}
    As a reminder, the ReLU activation function
    is $\mathrm{ReLU}(t) := \max(t, 0)$.
\end{remark}

To prove this theorem,
we'll need the following lemma.

\begin{theorem}{ReLU Breakpoints}
    If $h \colon \mathbb{R} \to \mathbb{R}$
    is a ReLU network with hidden layer widths
    $m_1, m_2, \dots, m_L$, then
    $h$ is a piecewise linear function
    with at most $2^{L - 1} (m_1 + 1)m_2 \dots m_L$
    linear pieces.
\end{theorem}

\begin{proof}
    Induction on $L$.  The base case of $L = 1$ states that
    a linear combination of $m$ ReLUs
    is piecewise linear with at most $m + 1$
    linear pieces, which is obviously true.

    For the inductive step, the point is the following:
    \begin{quote}
        \textbf{Claim.}
        Passing a piecewise linear function
        with $T$ linear pieces through a ReLU
        at most doubles the number of linear pieces.

        \emph{Justification:}
        Each linear piece
        either stays a single piece
        or gets split in two,
        depending on if it crosses the $x$-axis.
    \end{quote}
    Each pre-activation in the $L^{\text{th}}$ hidden layer
    is the output of a ReLU network with $L - 1$ layers
    and hidden layer widths $m_1, m_2, \dots, m_{L - 1}$.
    Thus, each $L^{\text{th}}$-layer pre-activation
    has at most $N := 2^{L - 2}(m_1 + 1)m_2\dots m_{L - 1}$ linear pieces
    by the inductive hypothesis.

    Then the claim above implies
    that each $L^{\text{th}}$-layer neuron has at most $2N$
    linear pieces.
    And the final output is a linear combination of $m_L$ such neurons,
    hence an upper bound of
    $2Nm_L = 2^{L - 1}(m_1 + 1)m_2\dots m_L$, as desired. ~ $\blacksquare$
\end{proof}

Now for the proof of the ``Depth Necessity'' theorem from earlier.

\begin{proof}
    Consider the tent map $\tau \colon [0, 1] \to [0, 1]$ given by:
    \[ \tau(x)
    := 2 \cdot \left[\mathrm{ReLU}(x) - \mathrm{ReLU}(2x - 1)\right]
    = \begin{cases}
    2x & x \leq \frac{1}{2}. \\
    2 - 2x & x > \frac{1}{2}.
    \end{cases} \]
    Then a depth-$O(L^2)$ network
    can produce the function $f(x) := \tau^{(\circ) L^2}(x)$
    (that is, $\tau$ composed with itself $L^2$ times).
    The result is a function with $2^{L^2 - 1}$ oscillations;
    here's how $f(x)$ looks for $L = 2$, for example.
    \begin{center}
        \frame{\includegraphics[width=0.3\textwidth]{lecture-13/iterated-tent}}
    \end{center}
    Any output $h \colon [0, 1] \to \mathbb{R}$
    of a ReLU network of depth $L$ and at most $2^L$ nodes
    has (according to the previous theorem):
    \[ \text{\# of linear pieces}
    \leq 2^{L - 1} \cdot \left[(m_1 + 1)m_2 \dots m_L\right]
    \leq 2^{L - 1} \cdot \left[\frac{2^L}{L}\right]^{L}
    \leq 2^{L^2 - 1}. \]
    We used some very crude estimates here,
    including an AM-GM bound.

    The point is that $f(x)$ forms $2^{L^2}$ triangles
    with the line $y = \frac{1}{2}$ (see the diagram above).
    But $h$ can only change direction at most $2^{L^2 - 1}$ times,
    so $h$ must be on the ``wrong side'' of at least
    half of the triangles.

    And the integral of $|f(x) - h(x)|$ over
    just the intervals where $h$ is on the ``wrong side''
    is already $\frac{1}{8}$, as desired. ~ $\blacksquare$
\end{proof}

\begin{remark}
    Some of our bounds here are really loose,
    enough so that the $\frac{1}{8}$ bound
    doesn't fully hold water as presented.
    But the holes are easy to fix if we
    settle for a smaller constant, say;
    the gist of the proof is still fully present.
\end{remark}

\subsection{Neural Network Optimization}

Neural networks are a parametric model
with parameters $\theta := \{W^{(i)}, b^{(i)}\}_{i = 1}^{L + 1}$.
As with any parametric model,
we tune the parameter $\theta$
by minimizing the empirical risk (ERM).
Specifically, we assume the following:
\begin{itemize}
    \item Our loss function is
    $\mathcal{L}\left(h_{\theta}(x^{(t)}), y^{(t)}\right) :=
    \left(h_{\theta}(x^{(t)}) - y^{(t)}\right)^2$
    (that is, squared loss).
    \item We employ a variant of (stochastic) gradient descent:
    for each training example $(x^{(t)}, y^{(t)})$, we perform the update:
    \[ \theta' \gets \theta - \alpha \cdot \nabla_{\theta} \left[
        \mathcal{L}\left(h_{\theta}(x^{(t)}), y^{(t)}\right)
        + \lambda \cdot \Omega(\theta) 
    \right]. \]
    Here, $\alpha$ and $\lambda$ are hyperparameters,
    and $\Omega(\theta) := \sum_{i = 1}^{L + 1} \|W^{(i)}\|^2_F$
    is a Frobenius regularizer.
    \begin{itemize}
        \item Note that $\|M\|^2_F$ denotes
        the sum of the squares of all of the entries in $M$.
        \item Recall from
        \href{/6.790/lectures/6/#map-for-bayesian-linear-regression}{Lecture 6}
        that regularizers (e.g. ridge regression terms)
        penalize large parameters
        to discourage overfitting.
    \end{itemize}
\end{itemize}

The gradient of the regularizer is easy to compute---the
gradient of $\mathcal{L}(h_{\theta}(x), y)$, less so.

\subsection{Chain Rule Warm-Up}

The rest of this lecture is dedicated to
the computation of the gradient of $\mathcal{L}(h_{\theta}(x), y)$.
First, a chain rule warm-up.

\begin{problem}
Write an expression for $\frac{\partial f}{\partial w}$
given the computational graph shown below.
\begin{center}
    \frame{\includegraphics[width=0.3\textwidth]{lecture-13/w-f-path}}
\end{center}
\end{problem}

\begin{solution}
    There is one term for each of the four paths from $w$ to $f$ in the graph;
    each term is the product of all of the partial derivatives
    along that path.
    \[
    \frac{\partial f}{\partial w}
    = \left[
        \frac{\partial f}{\partial h_1}
        \frac{\partial h_1}{\partial g_1}
        \frac{\partial g_1}{\partial w}
    \right]
    + \left[
        \frac{\partial f}{\partial h_1}
        \frac{\partial h_1}{\partial g_2}
        \frac{\partial g_2}{\partial w}
    \right]
    + \left[
        \frac{\partial f}{\partial h_2}
        \frac{\partial h_2}{\partial g_1}
        \frac{\partial g_1}{\partial w}
    \right]
    + \left[
        \frac{\partial f}{\partial h_2}
        \frac{\partial h_2}{\partial g_2}
        \frac{\partial g_2}{\partial w}
    \right].
    \]
    To help concretize this symbol slop,
    let's actually compute this partial derivative
    given the following functions:
    \[
    g_1(w) = 3w - 2
    ~~~ \parallel ~~~
    g_2(w) = w^2 + 1
    ~~~ \parallel ~~~
    h_1(g_1, g_2) = g_1g_2
    ~~~ \parallel ~~~
    h_2(g_1, g_2) = g_2^3 - g_1^2
    ~~~ \parallel ~~~
    f(h_1, h_2) = h_1^2h_2
    \]
    Let's compute $\frac{\partial f}{\partial w}$ when $w = 1$.
    The first step is a \emph{forward pass}:
    compute every function at $w = 1$.
    \begin{center}
        \frame{\includegraphics[width=0.6\textwidth]{lecture-13/chain-rule-example}}
    \end{center}
    Separately, we may compute all of the partial derivatives:
    \[
    \frac{\partial g_1}{\partial w} = 3
    ~~ \parallel ~~
    \frac{\partial g_2}{\partial w} = 2w
    ~~ \parallel ~~
    \frac{\partial h_1}{\partial g_1} = g_2
    ~~ \parallel ~~
    \frac{\partial h_1}{\partial g_2} = g_1
    ~~ \parallel ~~
    \frac{\partial h_2}{\partial g_1} = -2g_1
    ~~ \parallel ~~
    \frac{\partial h_2}{\partial g_2} = 3g_2^2
    ~~ \parallel ~~
    \frac{\partial f}{\partial h_1} = 2h_1h_2
    ~~ \parallel ~~
    \frac{\partial f}{\partial h_2} = h_1^2
    \]
    Notice that most of these expressions are \emph{not}
    direct functions of $w$.
    In order to evaluate them at $w = 1$,
    we must have done the precomputation
    of every function's value at $w = 1$
    via a \emph{forward pass}.

    Here's how it would look to compute
    the second of the four terms in $\frac{\partial f}{\partial w}$.
    \[
    \frac{\partial f}{\partial h_1} \times
    \frac{\partial h_1}{\partial g_2} \times
    \frac{\partial g_2}{\partial w}
    = 2(h_1)(h_2) \times g_1 \times 2w
    = 2(2)(7) \times (1) \times 2(1)
    = \boxed{56}.
    \]
    If we were to carry out
    the full computations for $\frac{\partial f}{\partial w}$,
    we would have to compute:
    \begin{itemize}
        \item All $6$ node values in the graph.
        (The \emph{forward pass}.)
        \item All $8$ partial derivative functions,
        and all $8$ of their evaluations at $w = 1$.
        \item A total of $8$ multiplications of partial derivatives
        and $3$ additions of products of partial derivatives.
    \end{itemize}
\end{solution}

\subsection{Neuron Warm-Up}

Again, as a warm-up,
we'll start by investigating the gradients
in the simplest possible neural network.

\textbf{Warm-Up.} In a single-neuron one-dimensional neural network,
here's how $\nabla_{\theta} \mathcal{L}(h_{\theta}(x), y)$ is computed.
\begin{center}
    \frame{\includegraphics[width=0.2\textwidth]{lecture-13/one-neuron-graph}}
\end{center}
Here, the pre-activation is $a \colon \mathbb{R} \to \mathbb{R}$
via $a(x) = b + wx$,
and the activation is $h(x) := g(a(x))$.
We follow the same approach as before;
first calculate each node's value via a forward pass,
then compute:
\[
\frac{\partial \mathcal{L}}{\partial w}
= \frac{\partial \mathcal{L}}{\partial h}
\cdot \frac{\partial h}{\partial a}
\cdot \frac{\partial a}{\partial w}
= 2(h(x) - y) \cdot g'(a(x)) \cdot x
~~~ \parallel ~~~
\frac{\partial \mathcal{L}}{\partial b}
= \frac{\partial \mathcal{L}}{\partial h}
\cdot \frac{\partial h}{\partial a}
\cdot \frac{\partial a}{\partial b}
= 2(h(x) - y) \cdot g'(a(x)) \cdot 1
\]
Note that there is only one path
from $w$ to $\mathcal{L}$ (as well as from $b$ to $\mathcal{L}$),
so there are no sums yet.

\subsection{Backpropagation Example}

Now suppose we have several layers,
each with several neurons.
We don't draw out each of the parameters
$\{(W^{(i)}, b^{(i)})\}_{i = 1}^{L + 1}$
as input nodes to each pre-activation $a^{(i)}$, though.
Instead, the drawing below shows the following:
\begin{itemize}
    \item The only nodes are the $\mathbb{R}$-valued
    pre-activations $a^{(i)}$
    and their activations $h^{(i)}$.
    \item Each pre-activation and its corresponding activation
    is given a unique index, written as a purple subscript.
    \item For each entry of each matrix $W^{(i)}$,
    that entry represents the weight
    of a directed edge from an activation $h^{(i - 1)}_u$
    to a pre-activation $a^{(i)}_v$.
    That weight is hence denoted $w[u \to v]$.
    \item We ignore the biases $\{b^{(i)}\}_{i = 1}^{L + 1}$ for now,
    though they operate almost identically.
\end{itemize}
Importantly, each $w[u \to v]$ edge
is an input to the pre-activation it feeds into;
the variable $w[u \to v]$ just doesn't get a node.
\begin{center}
    \frame{\includegraphics[width=0.9\textwidth]{lecture-13/NN-labeled}}
\end{center}
\begin{remark}
    The diagram above makes a minor mistake---the
    superscripts should read $\{(0), (1), (2), (3)\}$,
    as the leftmost layer shown is the input layer,
    not a hidden layer.
\end{remark}
Our ultimate goal is to compute 
$\frac{\partial \mathcal{L}}{\partial w[u \to v]}$
for all edges $w[u \to v]$ in the network.
However, notice the following:
\[
\frac{\partial \mathcal{L}}{\partial w[u \to v]}
= \frac{\partial \mathcal{L}}{\partial a_v}
\cdot \frac{\partial a_v}{\partial w[u \to v]}
= \frac{\partial \mathcal{L}}{\partial a_v}
\cdot h_u.
\]
It's easy to compute all values of $h_u$ via a forward pass;
it remains to compute
$\delta[v] := \frac{\partial \mathcal{L}}{\partial a_v}$
for all $v$.

This is the \emph{backward pass},
done via an induction-flavored approach.
The ``base case'' is $\delta[8]$,
the pre-activation node closest to $\mathcal{L}$.
\begin{center}
    \frame{\includegraphics[width=0.9\textwidth]{lecture-13/backprop-1}}
\end{center}
The green arrows show the only path from $a_8$ to $\mathcal{L}$;
in general, the base case will only have one such path.
This yields:
\[ \delta[8]
= \frac{\partial \mathcal{L}}{\partial a_8}
= \frac{\partial \mathcal{L}}{\partial h_8}
\cdot \frac{\partial h_8}{\partial a_8}
= 2(h_8 - y) \cdot g'(a_8). \]
Now, if we wish to compute $\delta[6]$,
the trick is to use our newfound knowledge of $\delta[8]$.
\begin{center}
    \frame{\includegraphics[width=0.9\textwidth]{lecture-13/backprop-2}}
\end{center}
By propagating backwards,
we minimize computational redundancies
by computing earlier values of $\delta$ in terms of
later ones.
\[ \delta[6] = \frac{\partial \mathcal{L}}{\partial a_6}
= \frac{\partial \mathcal{L}}{\partial a_8}
\cdot \frac{\partial a_8}{\partial h_6}
\cdot \frac{\partial h_6}{\partial a_6}
= \delta[8] \cdot w[6 \to 8] \cdot g'(a_6). \]
We can compute $\delta[7]$ in the same way.
With that layer complete, we now move back a layer
and tackle $\delta[3]$.
\begin{center}
    \frame{\includegraphics[width=0.9\textwidth]{lecture-13/backprop-3}}
\end{center}
In this case, there are two paths
from $a_3$ to the $3^{\text{rd}}$ layer,
so our derivative has two terms.
\[ \delta[3] = \frac{\partial \mathcal{L}}{\partial a_3}
= \left[
    \frac{\partial \mathcal{L}}{\partial a_6}
    \cdot \frac{\partial a_6}{\partial h_3}
    \cdot \frac{\partial h_3}{\partial a_3}
\right] +
\left[
    \frac{\partial \mathcal{L}}{\partial a_7}
    \cdot \frac{\partial a_7}{\partial h_3}
    \cdot \frac{\partial h_3}{\partial a_3}
\right]
= g'(a_3) \cdot \left[
    \delta[6] \cdot w[3 \to 6]
    + \delta[7] \cdot w[3 \to 7]
\right]. \]
And we can compute $\delta[4]$ and $\delta[5]$ in the exact same way.

\subsection{Backpropagation Statement}

Using the notation established in the previous example,
we can now describe the backpropagation algorithm.

\textbf{Backpropagation.}
After reading in a training example $(x, y)$\dots
\begin{itemize}
    \item First, perform a forward pass:
    precompute $h_i(x)$ and $g'(a_i(x))$
    for every activation $h_i$ and pre-activation $a_i$
    in the NN.
    \item Next, the backward pass: compute
    $\delta[u] := \frac{\partial \mathcal{L}}{\partial a_u}$
    for every pre-activation $a_u$ as follows.
    \begin{itemize}
        \item Say the output layer
        has pre-activation $a_F$.
        Then $\delta[F] = 2(h_F(x) - y) \cdot g'(a_F)$.
        \item For the rest, compute $\delta[u]$ inductively
        via $\delta[u] = g'(a_u) \cdot
        \left[
            \sum_{v \colon u \to v} \delta[v] \cdot w[u \to v]
        \right]$.
    \end{itemize}
    \item Update the parameters
    $\{(W^{(i)}, b^{(i)})\}_{i = 1}^{L + 1}$ like so:
    \[ w[u \to v] \gets w[u \to v] -
    \alpha \cdot \delta[v] \cdot h_u(x)
    ~~~ \parallel ~~~
    b[v] \gets b[v] - \alpha \cdot \delta[v]. \]
\end{itemize}
(You can check for yourself that
the updating formula for $b[v]$
is also correct.)

\end{document}